\section{Akuisisi dan Pra-pemrosesan Data}
Data primer dalam penelitian ini diperoleh melalui pengunduhan harga penutupan harian (\textit{closing price}) dari seluruh emiten yang terdaftar dalam indeks saham LQ45 melalui antarmuka pemrograman aplikasi (API) Yahoo Finance secara otomatis. Rentang waktu pengambilan data mencakup periode yang cukup panjang dan representatif guna memastikan bahwa model dapat menangkap berbagai dinamika pasar, mulai dari fase \textit{bullish} hingga fase koreksi sistemik yang ekstrem. Data mentah yang telah berhasil diperoleh kemudian disimpan di dalam struktur basis data sementara untuk dilakukan proses sanitasi serta pembersihan dari nilai-nilai kosong atau anomali data yang timbul akibat aksi korporasi. Penggunaan indeks saham LQ45 dipilih secara sengaja karena memiliki tingkat likuiditas yang tinggi serta kapitalisasi pasar yang besar, sehingga hasil optimasi yang diperoleh memiliki relevansi praktis yang signifikan bagi para investor institusional di Indonesia.

Proses binerisasi status pasar dilakukan dengan memetakan nilai imbal hasil kuantitatif ke dalam dua keadaan basis kuantum yang merepresentasikan dinamika pergerakan harga aset secara diskrit dan terstruktur. Dalam model ini, status $|0\rangle$ diberikan pada aset yang menunjukkan imbal hasil positif atau cenderung naik, sementara status $|1\rangle$ diberikan pada aset dengan imbal hasil yang bersifat negatif atau tetap selama periode observasi. Pemetaan biner ini bertujuan untuk mereduksi kompleksitas informasi tanpa menghilangkan fitur fundamental mengenai arah pergerakan aset yang sangat krusial dalam pengambilan keputusan investasi strategis. Kumpulan status biner ini kemudian disusun ke dalam sebuah matriks status temporal yang akan digunakan sebagai input utama untuk menghitung probabilitas gabungan serta derajat interaksi strategis antar aset pada modul penelitian selanjutnya.
